\subsubsection{Design science in information systems research (lämplig rubrik?)}
\paragraph{}
\hspace{-11}För att besvara forskningsfrågorna har Hevners metod för design science research tillämpats~\cite{Hevner2004}.
Metoden går ut på att på systematiskt vis skapa en innovativ IT-artefakt som effektivt löser ett tydligt business-problem~\cite{Hevner2004}.
Artefakten utvärderas och förbättras iterativt under studiens gång för att uppnå optimalt resultat~\cite{Hevner2004}.
Det cykliska förhållningssättet kan innefatta artefaktens skapande, tillämpning eller utvärdering.
Utöver att leda till förbättring av befintliga verktyg, system eller processer leder metoden även till djupare förståelse för det problemområde artefakten tar sikte på att förbättra~\cite{Hevner2004}.
Metoden har ett tydligt fokus på praktisk problemlösning~\cite{Hevner2004}.
\paragraph{}
\hspace{-11}En nyckelfaktor med Hevners design science research är vikten av att den är både relevant och rigorös~\cite{Hevner2004}.
%Med relevans avses vikten av att det finns ett verkligt problem som adresseras.
Med relevans avses vikten av att adressera ett verkligt problem som faktiskt finns någonstans, exempelvis i en organisation.
Målet ska vara att ta fram en teknisk lösning som på ett betydelsefullt och effektivt sätt löser problemet~\cite{Hevner2004}.
Med rigorositet avses stringens i tillämpningens alla delar för artefakten.
Genom att tillämpa och noggrant följa befintliga teorier och metoder säkerställs rigorositet~\cite{Hevner2004}.
Detta gäller allt från forskningen till design, planering och själva framtagandet av lösningen.
På dessa två hörnstenar vilar Hevners metod.
\paragraph{}
\hspace{-11}En uppsättning riktlinjer listas som kan följas i samband med både forskningen och framtagandet av artefakten för att underlätta processen att uppnå relevans och rigorositet~\cite{Hevner2004}.
Dessa riktlinjer är snarare vägledning än något som måste följas till punkt och pricka.
Däremot kan de vara ett stöd i att en forskningsstudie bedrivs metodiskt och systematiskt samt att IT-artefakten effektivt löser det identifierade problemet.
Vidare kan riktlinjerna hjälpa författaren med framtagandet av IT-artefakten.